

El presente capítulo tiene como propósito establecer las políticas, procedimientos y responsabilidades para la creación, control, distribución, revisión, almacenamiento y disposición final de la documentación técnica y administrativa asociada al mantenimiento de los UAS de SKYFILMS.\\

El objetivo principal es garantizar que toda la información utilizada o generada durante las actividades de mantenimiento sea confiable, trazable y esté actualizada, asegurando la integridad de los registros y la conformidad con las exigencias de la Aerocivil.\\

Este sistema documental constituye un elemento fundamental para demostrar el cumplimiento técnico y normativo de SKYFILMS ante cualquier auditoría o inspección.

\section{Alcance}
El procedimiento aplica a todos los documentos técnicos, registros, formularios y reportes que se generen, utilicen o archiven dentro del Área de Mantenimiento de SKYFILMS, incluyendo:

\begin{itemize}
    \item Manuales técnicos y programas de mantenimiento.

    \item Registros de inspección, órdenes de trabajo y reportes técnicos.

    \item Certificados de aeronavegabilidad y formularios de control.

    \item Registros digitales y físicos de componentes, baterías y herramientas.

    \item Correspondencia técnica emitida o recibida relacionada con mantenimiento.
\end{itemize}

Aplica tanto para documentos emitidos internamente como para los recibidos de fabricantes, proveedores o la Aerocivil.

\section{Principios de la Gestión Documental}
El sistema de gestión documental de SKYFILMS se basa en los siguientes principios operativos, orientados a garantizar el control, la trazabilidad y la autenticidad de la información técnica:

\begin{enumerate}
    \item \textbf{Integridad:} Ningún documento podrá ser modificado o destruido sin la debida autorización del Técnico UAS.

    \item \textbf{Vigencia:} Se utilizará exclusivamente la versión más actualizada de cada documento.

    \item \textbf{Trazabilidad:} Todo registro deberá permitir identificar su origen, fecha, responsable y propósito.

    \item \textbf{Accesibilidad:} El personal técnico tendrá acceso rápido y seguro a la información vigente necesaria para el desempeño de sus funciones.

    \item \textbf{Seguridad y confidencialidad:} La documentación deberá mantenerse protegida contra pérdida, daño, alteración o acceso no autorizado.
\end{enumerate}

Estos principios garantizan que la documentación técnica respalde con fidelidad la trazabilidad del mantenimiento y la aeronavegabilidad de los UAS.

\section{Tipos de Documentos}
La gestión documental de SKYFILMS abarca los siguientes tipos de documentos:

\begin{enumerate}
    \item \textbf{Documentos controlados:} Incluyen los manuales, formatos, procedimientos técnicos y programas aprobados por la organización. Están sujetos a control de versiones y deben ser actualizados formalmente.

    \item \textbf{Documentos no controlados:} Son copias informativas o de referencia, utilizadas solo para consulta y no válidas como evidencia técnica.

    \item \textbf{Registros técnicos:} Son los documentos que evidencian la ejecución del mantenimiento, como órdenes de trabajo, bitácoras de inspección, reportes de pruebas o certificados de calibración.
\end{enumerate}

Cada documento debe tener un código único, fecha de emisión, número de revisión y responsable de su elaboración o aprobación.

\section{Control y Actualización de Documentos}
El control de versiones es esencial para asegurar el uso de información válida en todo momento. Por lo tanto, SKYFILMS implementa el siguiente procedimiento:

\begin{enumerate}
    \item \textbf{Creación y aprobación:} Todo nuevo documento técnico será elaborado por el personal competente, revisado y aprobado por el Tecnico UAS antes de su emisión.

    \item \textbf{Numeración y codificación:} Los documentos se numeran y codifican según el sistema interno establecido, garantizando una estructura ordenada y uniforme.

    \item \textbf{Control de revisiones:} Toda modificación o actualización deberá registrarse en el historial de revisiones del documento, indicando el motivo, la fecha y el responsable del cambio.

    \item \textbf{Distribución:} La distribución se realizará únicamente por medios controlados (plataforma digital o copias físicas selladas), evitando la circulación de versiones no autorizadas.

    \item \textbf{Retiro de versiones obsoletas:} Las versiones reemplazadas se identificarán claramente como “NO VIGENTES” y se archivarán en la sección de documentos históricos, asegurando que no se usen operativamente.
\end{enumerate}

Este proceso garantiza que todos los documentos técnicos utilizados reflejen la información más reciente y aprobada por la organización.

\section{Almacenamiento y Protección de Documentos}
Los documentos físicos y digitales deben conservarse en condiciones que aseguren su protección y disponibilidad.

\begin{enumerate}
    \item \textbf{Archivo físico:} Los documentos impresos se almacenarán en el Archivo Técnico de Mantenimiento, organizado por tipo de documento y fecha de emisión. El acceso estará restringido al personal autorizado.

    \item \textbf{Archivo digital:} SKYFILMS dispone de un repositorio electrónico en la red interna, con respaldos automáticos semanales y control de acceso mediante contraseñas personales.

    \item \textbf{Conservación:} Los documentos deben mantenerse legibles, protegidos de humedad, calor o deterioro físico.

    \item \textbf{Respaldo:} Los archivos digitales se respaldan periódicamente en servidores seguros o medios externos, garantizando su recuperación en caso de pérdida o daño.
\end{enumerate}

Estas medidas aseguran que la documentación se mantenga disponible y protegida durante toda su vida útil.

\section{Tiempos de Retención y Disposición Final}
De acuerdo con los lineamientos de la Aerocivil y las políticas de SKYFILMS, los tiempos mínimos de conservación documental son los siguientes:

\begin{table}[h!]
\centering
\begin{tabular}{|c|c|}
\hline
\textbf{Tipo de Documento} & \textbf{Tiempo de Retención Mínimo} \\ \hline
Registros de mantenimiento y aeronavegabilidad & 5 años después de la última actualización \\ \hline
Certificados, inspecciones y bitácoras técnicas & 5 años o la vida útil del equipo, lo que sea mayor \\ \hline
Programas de mantenimiento y manuales técnicos & Vigentes mientras esté en uso \\ \hline
Registros administrativos o de capacitación técnica & 5 años \\ \hline

\end{tabular}
\caption{Retención y disposición final de documentos}
\label{tab:retencion}
\end{table}

Cumplido el período de retención, los documentos podrán ser archivados o destruidos según el procedimiento interno, siempre bajo autorización escrita del Tecnico UAS.

\section{Control Documental de Registros del MCM}

El Manual de Control de Mantenimiento (MCM) y todos sus anexos forman parte del sistema documental controlado de SKYFILMS.

Cualquier cambio, actualización o enmienda deberá ser:

\begin{itemize}
    \item Propuesta por el Técnico UAS.

    \item Aprobada formalmente por la Gerencia Responsable.

    \item Reportada a la Aerocivil, cuando se trate de una modificación sustancial del contenido aprobado originalmente.
\end{itemize}

Se mantendrá un registro de control de revisiones del MCM que documente cada \textbf{actualización, fecha de aprobación, número de revisión y descripción del cambio}.

\section{Responsabilidades}

Las responsabilidades específicas dentro del sistema de gestión documental se distribuyen de la siguiente manera:

\begin{enumerate}
    \item \textbf{Gerente General:} Establece las políticas generales de gestión documental y garantiza los recursos para su implementación.

    \item \textbf{Técnico de UAS:}
        \begin{itemize}
            \item Supervisa el cumplimiento del sistema documental, aprueba revisiones, controla la vigencia de los registros y autoriza la disposición final de documentos.


            \item Administra el archivo físico y digital, mantiene los registros actualizados, distribuye versiones controladas y verifica la correcta aplicación del sistema.

            \item Diligencian correctamente los formatos técnicos, garantizan la precisión de la información registrada y reportan cualquier pérdida o daño documental.
        \end{itemize}

\end{enumerate}

\section{Disposiciones Finales}
El cumplimiento estricto del presente procedimiento de Gestión Documental es obligatorio para todas las áreas técnicas y administrativas de SKYFILMS.\\

El uso de documentación no controlada, versiones obsoletas o registros incompletos será considerado una no conformidad técnica grave, sujeta a acciones correctivas inmediatas.\\

Una gestión documental adecuada permite asegurar la trazabilidad técnica, la transparencia operativa y la conformidad regulatoria de la organización, garantizando que todas las operaciones de mantenimiento se ejecuten bajo los más altos estándares de calidad y seguridad establecidos por la Aerocivil.